\section{Procedimientos de recolección de información}
\label{sec:procedimientos}

El proceso de generación, captura y registro de datos se ejecuta mediante un
procedimiento sistemático y secuencial estructurado en cinco fases operativas.
El procedimiento es determinista y replicable: su ejecución en cualquier equipo
con el mismo entorno produce resultados idénticos.

\subsection{Fase 1: Configuración y control de la semilla pseudoaleatoria}

Se fija de manera determinista la semilla maestra del generador
pseudoaleatorio en el valor \num{42}. Se inicializan el paso de discretización
$\Delta t = \SI{1,0}{\second}$, las matrices dinámicas $\mathbf{A}$,
$\mathbf{B}$ y $\mathbf{H}$, las covarianzas $\mathbf{Q}$ y $\mathbf{R}$, y las
condiciones iniciales del proceso y del estimador según la
tabla~\ref{tab:parametros}.

Las $M$ corridas independientes no reutilizan la misma secuencia aleatoria: cada
una recibe un flujo derivado de la semilla maestra mediante generación de
secuencias hijas, procedimiento que garantiza simultáneamente independencia
estadística entre corridas y reproducibilidad exacta del conjunto. Antes de
iniciar la simulación se verifica computacionalmente el rango de la matriz de
observabilidad, y la ejecución se aborta si la condición no se satisface.

\subsection{Fase 2: Ejecución del régimen operacional dinámico}

En cada iteración temporal $k \in [1, N]$ se inyecta la señal de control
$u_{k}$ correspondiente al perfil de descarga de camiones —cuatro ventanas de
\SI{45}{\second} a \SI{4,5}{\metre\cubed\per\second}, iniciadas en los pasos
\num{10}, \num{255}, \num{505} y \num{755}— y se propaga el estado real
$\mathbf{x}_{k}$ incorporando la perturbación de proceso $\mathbf{w}_{k}$.

\subsection{Fase 3: Degradación telemétrica}

Sobre la componente observable del estado real se aplica la inyección del ruido
estocástico de medición $v_{k}$, generando la lectura degradada $y_{k}$ que
alimenta simultáneamente al SCADA simulado y al módulo Filtro de Kalman. Esta
separación es metodológicamente esencial: el filtro nunca accede al estado real,
solo a la señal corrompida, en las mismas condiciones en que operaría sobre una
planta.

\subsection{Fase 4: Estimación recursiva}

El módulo ejecuta el ciclo predictor-corrector sobre la serie degradada. En
paralelo se calcula la línea base de la variable latente mediante diferencias
finitas de la señal cruda,
$\dot{V}^{\text{df}}_{k} = (y_{k} - y_{k-1})/\Delta t$, que constituye el único
procedimiento disponible para un SCADA convencional que quisiera estimar la tasa
de llenado sin un estimador de estados. Este término de comparación se incorpora
porque permite dimensionar de forma honesta el aporte del módulo sobre la
variable no sensada.

\subsection{Fase 5: Almacenamiento y serialización}

Los vectores de series de tiempo resultantes se consolidan y exportan a
registros CSV estructurados, junto con las métricas por corrida, los resultados
del análisis de sensibilidad y un archivo JSON de configuración y entorno. La
tabla~\ref{tab:evidencia} enumera la evidencia documental generada.

\begin{table}[H]
\centering
\caption{Registros de trazabilidad generados por el procedimiento.}
\label{tab:evidencia}
\footnotesize
\begin{tabular}{@{}p{5.0cm}p{3.0cm}p{5.4cm}@{}}
\toprule
\textbf{Archivo} & \textbf{Dimensión} & \textbf{Contenido}\\
\midrule
\texttt{serie\_temporal\_corrida01.csv} & 16 columnas $\times$ 1000 filas &
Series completas de estado real, lectura degradada, estimación, errores,
ganancia, covarianza e innovación\\
\addlinespace
\texttt{metricas\_montecarlo.csv} & 7 columnas $\times$ 50 filas &
Métricas de error de cada corrida independiente\\
\addlinespace
\texttt{barrido\_degradacion.csv} & 8 columnas $\times$ 7 filas &
Desempeño ante niveles crecientes de ruido real del sensor\\
\addlinespace
\texttt{malla\_desajuste\_QR.csv} & Malla $5\times5$ &
RMSE ante desajuste intencional de las covarianzas del filtro\\
\addlinespace
\texttt{resultados.json} & --- &
Parametrización completa, versiones del entorno y resumen consolidado\\
\bottomrule
\end{tabular}
\end{table}

\subsection{Resguardo de la consistencia del proceso}

La consistencia del procedimiento se resguarda por tres mecanismos. El primero
es la \textbf{ejecución desde un punto de entrada único}: un solo archivo
orquesta todas las fases, de modo que no existe la posibilidad de ejecutar
parcialmente el experimento con parámetros desactualizados. El segundo es la
\textbf{parametrización centralizada}: todos los valores del experimento residen
en una única estructura inmutable, y ningún módulo define constantes propias. El
tercero es la \textbf{verificación de reproducibilidad}, descrita en la
sección~\ref{sec:calidad}, que compara mediante suma de verificación las salidas
de ejecuciones sucesivas.
